\documentclass{article}
\usepackage{natbib}
\usepackage{graphicx}

\title{Complex QA and Language Models Hybrid Architectures Survey}
\author{Your Name}
\date{\today}

\begin{document}

\maketitle

\begin{abstract}
    This paper reviews the state-of-the-art of language model architectures and strategies for complex question-answering (QA), focusing on hybridization. Large Language Models (LLMs) excel in leveraging public data for standard problems but face challenges in handling specific, complex questions. Recent advancements and projects like ChatGPT and GALACTICA highlight both potential and limitations of LLMs in complex QA. We analyze required skills, evaluation techniques, challenges, current solutions, and emerging research trends in this domain.
\end{abstract}

\section{Introduction}
Complex question-answering (QA) tasks often require deep contextual understanding and nuanced responses beyond simple factual retrieval. While Large Language Models (LLMs) have shown remarkable capabilities in answering standard QA prompts, addressing complex queries such as those involving cultural variations or multi-dimensional problem-solving requires advanced architectures and strategies. This paper surveys the landscape of LLM architectures tailored for complex QA scenarios, emphasizing hybrid approaches that combine the strengths of LLMs with specialized knowledge and reasoning techniques.

\section{Required Skills and Evaluation Techniques}
Handling complex QA tasks demands a diverse skill set including domain knowledge, linguistic nuance comprehension, and contextual reasoning. Evaluation techniques for complex QA go beyond simple accuracy metrics, encompassing measures of fairness, robustness, toxicity, and overall reliability in diverse contexts.

\section{Challenges in Complex QA}
The complexity of QA tasks spans several challenges:
\begin{itemize}
    \item \textbf{Domain Adaptation}: Adapting LLMs to different knowledge domains and specific contexts.
    \item \textbf{Decomposition and Multi-step QA}: Breaking down complex queries into manageable sub-tasks and integrating multi-step reasoning.
    \item \textbf{Long-form and Non-factoid QA}: Addressing queries that require detailed explanations or opinions rather than factual answers.
    \item \textbf{Data Protection and Safety}: Ensuring sensitive data protection and mitigating risks of biased or harmful outputs.
    \item \textbf{Explainability and Truthfulness}: Providing transparent reasoning processes and ensuring outputs align with factual accuracy.
    \item \textbf{Temporal Reasoning}: Incorporating temporal aspects into QA processes for dynamic information retrieval and reasoning.
\end{itemize}

\section{Review of Existing Research}
This section synthesizes findings from community-edited research papers such as BIG, BLOOM, and HELM, which analyze the capabilities and limitations of LLMs in complex QA tasks. These studies provide benchmarks and insights into task complexity and evaluation metrics crucial for advancing research in this field.

\section{Current Solutions and Research Trends}
Promising approaches and trends in addressing complex QA challenges include:
\begin{itemize}
    \item \textbf{Hybrid LLM Architectures}: Integrating LLMs with structured knowledge bases and symbolic reasoning systems.
    \item \textbf{Training and Prompting Strategies}: Optimizing training protocols and prompt designs for enhanced performance in complex scenarios.
    \item \textbf{Human Reinforcement Learning}: Incorporating human feedback and reinforcement learning mechanisms to improve model adaptability and accuracy.
    \item \textbf{Neuro-symbolic and Structured Knowledge Grounding}: Combining neural networks with symbolic AI techniques for deeper semantic understanding.
    \item \textbf{Program Synthesis and Iterated Decomposition}: Automating complex task decomposition and solution synthesis through iterative refinement processes.
\end{itemize}

\section{Conclusion}
In conclusion, while LLMs have demonstrated significant potential in addressing complex QA challenges, advancements in hybrid architectures and novel research directions are crucial for overcoming existing limitations. Future research should focus on enhancing model robustness, interpretability, and adaptability to diverse cultural and domain-specific contexts.

\bibliographystyle{plainnat}
\bibliography{prompt1_try3_35}

\end{document}
